\chapter{Contrastive Embeddings: CLAP and the Vibe Axes \emph{(outline)}}\label{ch:contrastive}
\chapterlede{One model, two encoders, one shared space: CLAP aligns audio
with language, which lets us name directions in sound-space using English
sentences. The vibe visualizations are driven by dot products against those
directions.}

\begin{incode}
Status: outline. To be written against \texttt{backfill\_clap.py} (audio
tower), \texttt{compute\_vibe.py} (text axes: bright/warm/dense/tense/%
euphoric/vast), and \texttt{microscope.py}'s \texttt{/api/textsearch}.
\end{incode}

Planned contents:
\begin{itemize}
  \item \textbf{The InfoNCE objective.} Batch of $(a_i, t_i)$ pairs;
  \[
  \mathcal{L} = -\frac{1}{B}\sum_i \log
  \frac{\exp(\ip{a_i, t_i}/\tau)}{\sum_j \exp(\ip{a_i, t_j}/\tau)} ;
  \]
  symmetric audio$\to$text and text$\to$audio terms; the learned temperature
  $\tau$; InfoNCE as a lower bound on mutual information (statement, proof
  sketch).
  \item \textbf{Geometry of the shared space.} Why both towers L2-normalize;
  alignment vs.\ uniformity decomposition of the loss; the modality gap
  (audio and text occupy offset cones) and why \emph{differences} of text
  embeddings partially cancel it.
  \item \textbf{Zero-shot classification = nearest text.} Text search over
  moments as a single matrix--vector product (\texttt{emb @ text\_vec}).
  \item \textbf{Semantic axes.} The vibe construction: multi-prompt pole
  averaging (variance reduction), score $= \ip{e, v_+} - \ip{e, v_-}$;
  percentile normalization as a monotone calibration; why axes are
  \emph{named} but not \emph{orthogonal} --- correlated axes (bright/euphoric)
  and what that implies for palette mappings.
  \item \textbf{Failure modes.} Prompt sensitivity; the broken-checkpoint
  incident (\texttt{laion/larger\_clap\_music} collapsed embeddings) as a
  case study in validating learned components empirically.
\end{itemize}
